\begin{theorem}[A.4.7; Proved; Euler number and fundamental group for every toroidal output]\label{push:A47}
Let $\Sigma$ be any output of A.2.4 for $u=r^n$, and let
$D_\Sigma$ denote the underlying reduced central space.
Let $k(\Sigma)$ be the number of selected pure top ray
orbits, and let $t_4(\Sigma)$ be the number of maximal
four-cones in one fundamental prism. Then
\[
 \chi(D_\Sigma)=11k(\Sigma)+t_4(\Sigma),
 \qquad \pi_1(D_\Sigma)=\mathbb Z,
 \qquad H_1(D_\Sigma;\mathbb Z)=\mathbb Z.
\]
These assertions include nonreduced central divisors.
At a point with precisely one local central branch of
multiplicity $\ell$, the nearby Euler function is $\ell$;
at a point with at least two local branches it is zero.
Its integral is always two. For every reduced output the
pure top count is $k(\Sigma)=n$, so
$\chi(D_\Sigma)=11n+t_4(\Sigma)$.
The higher integral homology of every arbitrary output is
not asserted by this theorem.
\end{theorem}
\begin{proof}
Over $B\setminus\{\infty\}$ the top fan gives a cycle
of $k(\Sigma)$ ruled bundles over
$S\setminus S_\infty$, glued along their zero and
infinity sections with the specified translations.
There are no triple intersections there. Its Euler
number is therefore
\[
 2k(\Sigma)\chi(S\setminus S_\infty)
       -k(\Sigma)\chi(S\setminus S_\infty)
 =11k(\Sigma),
\]
because $\chi(S)=12$ and $\chi(S_\infty)=1$.
This computation concerns the support of the divisor,
so the multiplicities do not enter.

At $q=r=0$ the strata are the torus orbits of cones whose
generators collectively have positive $q$-height and
positive $r$-height. Each such cone has trivial period
stabilizer. Indeed, an element $G_1^aG_2^b$ preserving
its finite set of primitive rays has a power fixing them
all. A ray of positive $q$-height then forces $a=0$;
any ray with positive $q$- or $r$-height forces $b=0$.
Consequently a stratum of cone dimension $d$ is an actual
$(\mathbb C^*)^{4-d}$, with no additional torus quotient.
The induced gluing units only identify different period
translates of the stratum. Its compactly supported Euler
number is zero unless $d=4$, in which case it is one.
Every maximal cone has both kinds of height, and each has
one representative in the interior of the fundamental
prism. Thus the cusp contributes exactly $t_4(\Sigma)$.
Additivity proves the displayed central Euler formula.

The local total space is smooth, and its parameter is
$r=z_1^{\ell_1}\cdots z_s^{\ell_s}$ times a holomorphic
unit, with $\ell_i>0$. Its Milnor fibre has the homotopy
type of a finite disjoint union of $(S^1)^{s-1}$.
For $s=1$ it consists of $\ell_1$ contractible pieces;
for $s\geq2$ its Euler number is zero. This proves the
pointwise assertion. Off the cusp the one-branch parts
are $\mathbb C^*$-bundles over the appropriate parts of
$S\setminus S_\infty$ and have Euler number zero; the
double parts have nearby Euler value zero. At the cusp,
positive-dimensional torus strata again contribute zero.
A maximal cone with exactly one positive-$r$ ray must
have three bottom rays. Its bottom face is consequently
one of the two prescribed triangles $ABC$ and $BCD$.
Each triangle belongs to exactly one such maximal cone.
Its determinant is the $r$-height $\ell$ of the fourth
ray, because the bottom triangle is unimodular.
Smoothness forces $\ell=1$. Exactly these two points
contribute one each. Thus the nearby Euler integral is
two for every output, even when some other components
have multiplicity greater than one.

For the fundamental group, first choose the reference
model obtained by base-changing $\mathcal X^+$ by
$u=r^n$ and resolving the surface equation $xy=r^n$
along its smooth conductor $S$. Its normalized central
components are $R^+$ and $n-1$ projective-line bundles
over $S$. All are simply connected, and their double
surfaces are copies of the simply connected $S$.
The gluing graph is an $n$-cycle, interpreted as the
self-gluing circle when $n=1$. Van Kampen therefore
gives fundamental group $\mathbb Z$ for this reference
centre.

Any two fans of A.2.4 have the explicit proper common
normal modification formed by intersecting their cones.
Resolve this analytic space to a smooth common model;
analytic resolution, including the noncompact version
near the compact central set, is supplied by
\cite[Theorem 1.6 and its following paragraph]{BM95AB}.
This use of a common modification is legitimate here
because the common fan has already constructed it.
It is not an assumed extension of an arbitrary punctured
isomorphism.

A proper bimeromorphic map $F:Y\to X$ between smooth
complex manifolds induces an isomorphism of fundamental
groups. To see this directly, choose an analytic
$A\subset X$ of complex codimension at least two such
that $F$ is an isomorphism outside $A$. General position
gives $\pi_1(X\setminus A)\simeq\pi_1(X)$, and gives
surjectivity of
$\pi_1(Y\setminus F^{-1}A)\to\pi_1(Y)$.
Under the common-open identification, composition with
$F_*$ is the displayed isomorphism to $\pi_1(X)$.
Thus the first map is also injective, and $F_*$ is an
isomorphism. Apply this argument to both maps from the
smooth common model, after restricting to a sufficiently
small parameter disc.

Finally, each smooth total space retracts to the support
of its central divisor, including in the nonreduced case.
The gradient proof in B.4.2 applies with the explicit
bound for $\rho=|r|^2$
\[
 \|\nabla\rho\|\geq c\rho^{1-1/(2L)},
 \qquad L=\ell_1+\cdots+\ell_s.
\]
On the complement of the divisor it follows from the
weighted arithmetic--geometric mean inequality applied
to $\sum_i\ell_i^2/|z_i|^2$. A holomorphic unit can be
absorbed into one coordinate locally. Compactness of
the central set gives a uniform exponent below one,
and the resulting finite remaining path-length bound
gives the continuous gradient retraction. Hence all
the centres and their smooth total neighbourhoods have
the same fundamental group as the reference centre.
Abelianization gives the stated $H_1$.

The program \path{code/ab_push_all_fan_topology.py}
checks the Euler and two-point nearby formulas for every
minimal bridge at $n=1,2,12$, eight point-star
modifications, and five nonreduced conductor-star
modifications. In the latter family the added pure top
rays are $(0,1,0,j)$ for $2\leq j\leq m+1$; the formula
gives $\chi(D_m)=14+12m$, independently agreeing with
successive blow-ups of the conductor surface of Euler
number twelve. The finite checks are recorded separately
from the proof for every fan.
\end{proof}
